% Generated draft: requires mathematical and source review.
\documentclass[11pt]{article}
\usepackage[margin=25mm]{geometry}
\usepackage{fontspec}
\setmainfont{Latin Modern Roman}
\usepackage{amsmath,amssymb,mathtools,unicode-math}
\setmathfont{Latin Modern Math}
\usepackage{enumitem,needspace,xurl,hyperref}
\setlength{\parindent}{0pt}
\setlength{\parskip}{5pt}
\setlength{\emergencystretch}{3em}
\newcommand{\N}{\mathbb{N}}
\newcommand{\Z}{\mathbb{Z}}
\newcommand{\Q}{\mathbb{Q}}
\newcommand{\R}{\mathbb{R}}
\newcommand{\C}{\mathbb{C}}
\newcommand{\F}{\mathbb{F}}
\newcommand{\A}{\mathbb{A}}
\newcommand{\PP}{\mathbb{P}}
\newcommand{\E}{\mathbb{E}}
\newcommand{\eps}{\varepsilon}
\newcommand{\dd}{\,\mathrm{d}}
\begin{document}
\section*{UnsolvedMath ID: 800071}
\textbf{Source problem code:} catalogue entry
\subsection{Definitions and mathematical statement}
A contact structure on a smooth three-dimensional manifold \(M\) is a rank-two subbundle \(\xi\subset TM\) that locally has the form \(\xi=\ker\alpha\), where \(\alpha\) is a smooth one-form satisfying \(\alpha\wedge\dd\alpha\neq 0\) at every point of its domain. A line field tangent to \(\xi\) is a smooth rank-one subbundle \(L\subset\xi\); this does not require a chosen orientation or a global nonzero section of \(L\).

A smooth two-form \(\omega\in\Omega^2(M)\) is closed and nowhere vanishing if
\[
\dd\omega=0,\qquad \omega_p\neq 0\quad\text{for every }p\in M.
\]
It annihilates \(L\) if
\[
\omega_p(v,w)=0\quad\text{for every }p\in M,\ v\in L_p,\ w\in T_pM.
\]
In dimension three, these conditions imply \(\ker\omega_p=L_p\). Interpret ``of integral class'' as
\[
[\omega]_{\mathrm{dR}}\in\operatorname{im}\bigl(H^2(M;\Z)\longrightarrow H^2_{\mathrm{dR}}(M;\R)\bigr),
\]
where the map is induced by change of coefficients and the de Rham isomorphism. The supplied record states no alternative normalization.

\textbf{Question (Nariman--Yazdi, Question 12).} How can one systematically construct examples of triples \(((M,\xi),L,\omega)\) satisfying all the conditions above?

The request concerns a systematic production of examples; the source does not assert existence for every prescribed contact manifold or every prescribed tangent line field. It does not specify compactness, connectedness, boundary conditions, or coorientability of \(\xi\), and it does not define what qualifies as a systematic construction. These qualifications remain unresolved rather than being added as hypotheses.
\subsection{Short English statement}
Find systematic constructions of contact three-manifolds equipped with a line field inside the contact planes and a nowhere-vanishing closed two-form of integral cohomology class whose kernel is that line field. The source leaves the intended scope of the examples and the meaning of ``systematic'' unspecified.
\subsection{Sources}
\begin{itemize}
\item Local UnsolvedMath registry: permanent ID \(800071\), exact source problem code \texttt{catalogue entry}, title ``Foliations and Diffeomorphism Groups --- Question 12.'' The supplied record attributes the question to Nariman and Yazdi, with proposed year 2025, and labels its status \texttt{open}; this status is supplied registry data, not an independently verified mathematical assessment.
\item Nariman and Yazdi, \emph{Problem List on Foliations and Diffeomorphism Groups} (2025), Question 12. Supplied source URL: \url{https://arxiv.org/abs/2503.18274}. The supplied provenance records access on 2026-10-16 and extraction by \texttt{arxiv\_latex}; these are supplied provenance details, not a claim of independent source access or verification.
\end{itemize}
\end{document}
